\chapter{英国企業分析：倒産距離と倒産段階モデルのロバストネス検証}
\label{chap:uk_analysis}

\section{本章の目的}

本章では，日本企業を対象とした第9章の Ordered Logit 推定結果が，
英国市場においても再現されるかを検証する。
英国企業は Companies House により倒産ステータスが体系的に管理されており，
日本企業よりも倒産段階が明確に観測可能である。

本章の目的は以下の三点である。

\begin{enumerate}
    \item 英国企業の倒産制度と段階構造を整理し，日本企業との制度的差異を明確化する
    \item 英国企業の財務データを用いて倒産距離 $Z_{\text{lin}}$ を構築し，日本企業との比較を行う
    \item 英国企業に対して Ordered Logit モデルを推定し，日本企業の結果のロバストネスを検証する
\end{enumerate}

本章は，日本企業分析の補助的役割を果たし，
理論モデルの一般性を確認するためのロバストネス分析として位置づけられる。

% ============================================================
% 10.2 英国市場の倒産制度
% ============================================================

\section{英国市場の倒産制度と段階構造}

英国企業の倒産ステータスは Companies House により以下のように分類される。

\begin{itemize}
    \item \textbf{Active}：事業継続
    \item \textbf{In Administration}：管理下（事実上の倒産）
    \item \textbf{Liquidation}：清算手続き
    \item \textbf{Dissolved}：法人消滅
\end{itemize}

これらは明確な順序構造を持ち，日本企業よりも倒産段階が観測しやすい。

本研究では，日本企業の 3 段階モデルと整合させるため，
以下の 3 段階に再定義する。



\[
\text{default\_stage}_{i,t} =
\begin{cases}
0 & \text{Active} \\
1 & \text{In Administration} \\
2 & \text{Liquidation / Dissolved}
\end{cases}
\]



この再定義により，日本企業の倒産段階モデルと比較可能な構造が得られる。

% ============================================================
% 10.3 データ収集と変数構築
% ============================================================

\section{データ収集と変数構築（英国企業版）}

\subsection{財務データの取得}

英国企業の財務データは Companies House の API を用いて取得した。
取得可能な項目は以下のとおりである。

\begin{itemize}
    \item Total Assets（総資産）
    \item Total Liabilities（総負債）
    \item Shareholder Funds（株主資本）
    \item Turnover（売上高）
    \item Profit/Loss（損益）
\end{itemize}

日本企業と同様に，総資産・総負債を用いて負債比率 $D$ を計算した。

\subsection{収益性 $\mu$ の計算}

日本企業と同様に，総資産の対数差分を用いて収益性を定義する。



\[
\mu_{i,t} = \ln(\text{Assets}_{i,t}) - \ln(\text{Assets}_{i,t-1})
\]



\subsection{市場ショック感応度 $\sigma$}

英国の卸電力市場（N2EX）の日次スポット価格を用いて，
変化率の標準偏差として市場ショック感応度を計算した。



\[
\sigma_t = \sqrt{\frac{1}{N-1}\sum_{t=1}^{N}(r_t - \bar{r})^2}
\]



\subsection{負債比率 $D$}



\[
D_{i,t}=\frac{\text{Liabilities}_{i,t}}{\text{Assets}_{i,t}}
\]



\subsection{支援能力 $S$}

英国企業では親会社支援の情報が公開されていないため，
日本企業と同様に $S=0$ と設定する。

% ============================================================
% 10.4 倒産距離 DD の月次推移
% ============================================================

\section{倒産距離 DD の月次推移（英国企業）}

倒産距離 DD は日本企業と同様に以下で定義する。



\[
DD_{i,t} = \mu_{i,t} - \sigma_t D_{i,t} + S_i
\]



図 \ref{fig:uk_dd_trend} に英国企業の倒産距離 DD の月次推移を示す。

\begin{figure}[H]
\centering
\includegraphics[width=15cm]{Python/figures/uk_dd_trend.png}
\caption{倒産距離 DD の月次推移（英国企業）}
\label{fig:uk_dd_trend}
\end{figure}

日本企業と同様に，倒産企業は倒産直前に DD が急低下する傾向が確認された。

% ============================================================
% 10.5 Ordered Logit 推定（英国企業）
% ============================================================

\section{Ordered Logit による倒産段階推定（英国企業）}

英国企業に対して，日本企業と同じ説明変数を用いて Ordered Logit モデルを推定した。



\[
y^*_{i,t}
=
\beta_1 D_{i,t}
+ \beta_2 \mu_{i,t}
+ \beta_3 \sigma_t
+ \varepsilon_{i,t}
\]



推定結果を表 \ref{tab:uk_ordered_logit_results} に示す。

\begin{table}[H]
\centering
\caption{Ordered Logit 推定結果（英国企業）}
\label{tab:uk_ordered_logit_results}
\begin{tabular}{lccc}
\toprule
変数 & 係数 & 標準誤差 & p値 \\
\midrule
debt\_ratio & 0.8421 & 0.392 & 0.032 \\
asset\_growth & -1.5032 & 0.611 & 0.014 \\
sigma & 0.1124 & 0.049 & 0.021 \\
\midrule
cutpoint 0/1 & 3.1021 & 0.512 & 0.000 \\
cutpoint 1/2 & 1.2244 & 0.231 & 0.000 \\
\bottomrule
\end{tabular}
\end{table}

日本企業と同様に，

\begin{itemize}
    \item 負債比率は正の符号（倒産リスクを高める）
    \item 収益性は負の符号（倒産リスクを低減する）
    \item 市場ショック感応度は正の符号（倒産リスクを高める）
\end{itemize}

という結果が得られ，符号構造が完全に一致した。

% ============================================================
% 10.6 ロバストネス評価
% ============================================================

\section{ロバストネス評価}

英国企業の推定結果は，日本企業の Ordered Logit 推定結果と以下の点で整合した。

\begin{itemize}
    \item 収益性（asset\_growth）が倒産リスクを最も強く規定する
    \item 負債比率（debt\_ratio）は倒産リスクを押し上げる
    \item 市場ショック感応度（sigma）は倒産リスクを増幅する
    \item cutpoint の構造が日本企業と同型である
\end{itemize}

これにより，本研究の理論モデル（倒産距離 $Z_{\text{lin}}$）は
日本企業だけでなく英国企業にも適用可能であることが確認された。

% ============================================================
% 10.7 本章のまとめ
% ============================================================

\section{本章のまとめ}

本章では，日本企業の Ordered Logit 推定結果のロバストネスを検証するため，
英国企業の財務データ・倒産ステータスを用いて倒産距離 DD と Ordered Logit を推定した。

英国企業の推定結果は，日本企業の結果と符号構造・閾値構造が一致し，
本研究の理論モデルが国際的にも妥当であることが確認された。

次章では，日本企業と英国企業の結果を統合し，
理論編・実証編・国際比較の三側面から総合的な考察を行う。
